% Pandoc template for the fed_fraud documents.
% Based on proposal_gen/templates/detailed_proposal_template/detailed_1_classic.tex
% (Classic report): centered title page, table of contents, ruled section
% headings and the Geodesic house footer. Fonts are Latin Modern (lmodern).
% Two layouts: the full report (title page and contents) and, with
% -V brief, a compact one-page-brief header in the same style.
\documentclass[letterpaper,11pt]{article}
\usepackage[T1]{fontenc}
\usepackage[utf8]{inputenc}
\usepackage{lmodern}
\usepackage[letterpaper,top=0.9in,bottom=1.2in,left=1in,right=1in,footskip=0.6in]{geometry}
\usepackage[table]{xcolor}
\usepackage{geodesicbrand}
\usepackage{amsmath,amssymb}
\usepackage{enumitem}
\usepackage{array}
\usepackage{calc}
\usepackage{longtable}
\usepackage{makecell}
\usepackage[most]{tcolorbox}
\usepackage{titlesec}
\usepackage{needspace}
\usepackage{etoolbox}
\usepackage{hyperref}
\hypersetup{colorlinks=true,urlcolor=coolblack,linkcolor=coolblack,filecolor=geoblue,
  pdftitle={Energy},pdfauthor={\geofirm}}
\urlstyle{same}

\setlist{itemsep=2pt,leftmargin=1.4em}
\setlength{\parindent}{0pt}
\setlength{\parskip}{6pt}
\providecommand{\tightlist}{\setlength{\itemsep}{0pt}\setlength{\parskip}{0pt}}
\setcounter{secnumdepth}{0}

% Section style: house rule under each heading
\titleformat{\section}{\Large\bfseries\color{coolblack}}{\thesection}{0.6em}{}[{\color{coolblack}\titlerule[0.5pt]}]
\titleformat{\subsection}{\large\bfseries\color{geomid}}{\thesubsection}{0.6em}{}
\titleformat{\subsubsection}{\normalsize\bfseries\color{geomid}}{\thesubsubsection}{0.6em}{}
\titlespacing*{\section}{0pt}{10pt}{6pt}

% Tables are written by latex/filters.lua (navy header, banded rows)
\setlength{\LTpre}{6pt}\setlength{\LTpost}{6pt}

\begin{document}

% BRIEF HEADER
\thispagestyle{geodesic}
\begin{minipage}[c]{0.13\textwidth}\includegraphics[width=\linewidth]{icon.png}\end{minipage}\hfill
\begin{minipage}[c]{0.84\textwidth}
  {\textsc{\large \geofirm}}\par\vspace{2pt}
  {\small\color{geogray}Federal fraud expansion\enspace\textperiodcentered\enspace When: \compiletime}
\end{minipage}\par\vspace{6pt}
{\color{coolblack}\rule{\linewidth}{0.3mm}}\par\vspace{2pt}
{\LARGE\bfseries\color{coolblack} Energy}\par
\vspace{2pt}{\large\color{geomid} One-page brief}\par
\vspace{-2pt}{\color{coolblack}\rule{\linewidth}{0.3mm}}\par\vspace{4pt}
\pagestyle{geodesic}

\textbf{Priority:} 11 of 33 (score 3.3). \textbf{FY2026 outlays:} \$63B
(USAspending). \textbf{As of:} 2026-10-27. Public data only.

\hypertarget{bottom-line}{%
\section{Bottom line}\label{bottom-line}}

Most DOE money goes to national lab operating contractors and nuclear
cleanup, which are hard to score from outside. The fraud-prone slice is
the fast-growing grant, loan and small business research money from the
2021 infrastructure law and 2022 Inflation Reduction Act: many new
recipients, large awards, and thin track records. That is where public
award data and the SBA entity checks fit.

\hypertarget{exposure}{%
\section{Exposure}\label{exposure}}

\begin{itemize}
\tightlist
\item
  \$63B in outlays and \$141B in budgetary resources in FY2026, swollen
  by multi-year infrastructure and clean energy appropriations.
\item
  DOE is one of the largest SBIR/STTR funders. GAO found almost 8\% of
  SBIR and STTR awardees had four or more fraud risk indicators and
  about \$445M in possibly duplicative awards across agencies (2024).
\item
  DOE does not report a high-risk improper payment program.
\end{itemize}

\hypertarget{public-data-we-can-use}{%
\section{Public data we can use}\label{public-data-we-can-use}}

\begin{itemize}
\tightlist
\item
  \textbf{USAspending and FPDS:} grants, cooperative agreements,
  contracts and subawards by recipient.
\item
  \textbf{SBIR.gov award data:} firm, PI, abstract and agency for every
  award.
\item
  \textbf{DOE Loan Programs Office portfolio} listings.
\item
  \textbf{SAM.gov} entity and exclusion records; DOE OIG reports.
\end{itemize}

\hypertarget{detection-plays}{%
\section{Detection plays}\label{detection-plays}}

\begin{enumerate}
\def\labelenumi{\arabic{enumi}.}
\tightlist
\item
  \textbf{Duplicate research awards:} near-identical SBIR abstracts
  funded by DOE and other agencies (text similarity across SBIR.gov).
\item
  \textbf{Thin recipients, big awards:} new entities winning large
  grants, shared addresses across ``independent'' applicants, and
  pass-through subawards.
\item
  \textbf{Excluded links:} recipients or principal investigators tied to
  excluded parties or prior settlements.
\end{enumerate}

\hypertarget{first-prototype}{%
\section{First prototype}\label{first-prototype}}

A cross-agency SBIR duplicate-award screen covering DOE, NSF, NASA and
DoD. It is cheap (one public data set) and serves four briefs at once.

\hypertarget{limits}{%
\section{Limits}\label{limits}}

Lab and cleanup contractor costs are not visible below the award level.
Duplicate abstracts are leads; some overlap is legitimate.

\hypertarget{who-to-take-it-to}{%
\section{Who to take it to}\label{who-to-take-it-to}}

DOE OIG, DOE Office of Science SBIR/STTR program office, and the SBA
SBIR program office.

\hypertarget{sources}{%
\section{Sources}\label{sources}}

\begin{itemize}
\tightlist
\item
  \href{https://www.gao.gov/products/GAO-24-105470}{GAO-24-105470:
  SBIR/STTR fraud vulnerabilities (Sept 2024)}
\item
  \href{https://www.sbir.gov/awards}{SBIR.gov awards data}
\item
  \href{https://www.usaspending.gov/agency}{USAspending agency profiles
  (FY2026 outlays)}
\end{itemize}

\end{document}
